\section{Conclusion}
\label{sec:conclusion}

We asked whether a language model's internal state, partway through a familiar
multi-word expression, already encodes the rest. For fixed expressions in
natural text the answer is yes, and at a useful range: OLMo-2-7B completes such
phrases three tokens early in roughly two thirds of instances and five tokens
early in more than half. The ability is a stable property of the phrase rather
than of the occasion: across a phrase's natural contexts, success clusters at
the extremes far more than independent per-context variation would produce. It
is also legible from the representation itself, well enough that a small probe reading a single hidden
vector predicts real completion at $87.6\%$ precision.

The second finding is a caution about how the first was obtained. Patchscopes
and ordinary generation, run on identical instances, agree on only $78.7\%$ of
them, reverse the ranking of the phrase categories, and, used as probe labels,
locate the relevant information in different parts of the network. Part of that
gap is a budget we chose, since widening the sweep from eight layers to all $32$
recovers six to eight points, but the disagreement itself survives the widest
sweep we can run. Neither instrument is wrong; neither is a neutral window onto
what the model knows. Claims of the form \emph{the model already knows $X$}
should carry the readout that produced them, and the budget it was given.

The practical thread we have not pulled is efficiency. A signal that is stable
per phrase, available from mid-network states, and precise enough to act on is
the kind of signal a decoding scheme could use to commit several tokens at once
rather than paying full depth for each. We are pursuing that direction
separately, as an adaptive multi-token prediction policy checked token by token
against an autoregressive baseline; nothing in the present paper depends on its
outcome. An early pilot there is instructive about the size of
the prize: reproducing CLP \citep{xie-zhou-2026-clp}, a learned linear read of
the hidden state that sets the draft length per step, on a different model, we
measured roughly $1.34\times$ the throughput of plain autoregressive decoding.
A policy that simply drafts a fixed number of tokens at every step was faster
still, at $1.72\times$, while diverging from the autoregressive output
more often.\footnote{Three caveats on that pilot. Our figure exceeds the
$1.14\times$--$1.29\times$ Xie and Zhou report, which we take to reflect
differences in model, hardware and workload rather than a stronger
implementation. We observed occasional divergences from autoregressive output
under \emph{every} drafting policy we tried, including fixed-length ones with no
adaptive logic at all, which points at this model's interaction with the
inference engine rather than at any policy. And those divergences were more
common for the more aggressive fixed-length policies (six of ten prompts,
against two for CLP), so the fixed-length speed advantage is not free.}
Ten held-out prompts settle nothing either way. What the pilot does suggest is
that the open question is not whether anticipation exists, which this paper settles,
but whether reading it out adaptively buys enough over drafting blindly to pay
for itself.
